\barrerapagina
\section{Comparación académica entre Ventas y Transacciones}
La metodología comparte modelos multidimensionales propios, orden temporal, selección en validación y transformaciones ajustadas solo con train. Sin embargo, no son dos mediciones sobre la misma población: Ventas predice unidades relativas por producto provisional; Transacciones predice frecuencia relativa de cobros por cuenta operacional. Las columnas, identidad, escala empresarial, supuestos y baselines son específicos de cada contrato.

Ventas utiliza 20 características, con lags de unidades y variables históricas comerciales/monetarias; Transacciones conserva 16, centradas en conteos, estadísticos, antigüedad observada y calendario. En esta última no se habilitan importes ni saldos históricos no acreditados. Las ventanas principales de test contienen 1\,111 y 1\,284 observaciones, respectivamente; sus positivos no tienen el mismo significado.
\begin{table}[htbp]
\centering\small
\caption{Contraste descriptivo entre problemas distintos: porcentajes de test, sin comparación estadística emparejada.}\label{tab:tf-comparacion}
\begin{tabular}{llrrrrr}\toprule
Proceso & Modelo & N & Acc. & BA & Recall & Espec. \\ \midrule
Ventas & Perceptrón & 1111 & 53,915 & 54,104 & 96,022 & 12,186 \\
Ventas & Logística & 1111 & 62,286 & 62,415 & 91,139 & 33,692 \\
Transacciones & Perceptrón & 1284 & 62,461 & 66,222 & 86,857 & 45,586 \\
Transacciones & Logística & 1284 & 65,265 & 62,809 & 49,333 & 76,285 \\
\bottomrule\end{tabular}
\end{table}

La tabla \ref{tab:tf-comparacion} conserva cada población. La mayor BA observada en Ventas corresponde a logística (62,42\%); en Transacciones, al perceptrón (66,22\%). No procede restar ambas cifras como estimación de superioridad sobre un mismo problema ni concluir que un algoritmo sea universalmente mejor. Se comparan dos experimentos con diferentes etiquetas y muestras, no un ensayo emparejado de métodos en una única población.

En Ventas, ambos modelos presentan alta sensibilidad y escasa especificidad, particularmente el perceptrón. La logística reduce falsas alertas y obtiene mejor equilibrio, aunque sus 370 FP todavía limitan el uso automático. En Transacciones, el perceptrón reconoce muchos positivos con FPR 54,41\%, mientras logística sacrifica recall para obtener mayor especificidad y accuracy. La interpretación debe partir del costo y la acción empresarial, que no se han cuantificado; ninguna métrica equivale a rentabilidad o conciliación bancaria.

El perceptrón continúa actualizando en ambos procesos, pero sus checkpoints elegidos son 70 y 3. La logística tiene descenso estable en ambos: Ventas agota 200 épocas con tolerancia desactivada; Transacciones agota 400 sin alcanzar $10^{-6}$. No se interpreta la distinta época elegida como evidencia de convergencia más rápida a una solución comparable. Los PCA retienen 49,619\% y 45,021\% de varianza; solo visualizan estructuras distintas y no intervienen en las predicciones.

Las dos pruebas cubren cinco meses y tienen dependencia temporal y por entidad. La selección repetida en validación introduce optimismo de desarrollo; los pocos productos/cuentas nuevos impiden extrapolaciones estables. Los porcentajes son descriptivos de WWI y sus supuestos de cobertura. No se han demostrado significación estadística, deriva conceptual causal, superioridad universal ni utilidad productiva.

\section{Conclusiones integradas y continuidad del alcance académico}
La aplicación de Ventas conserva sus conclusiones verificadas: preparación causal común de 20 entradas, comparación reproducible de modelos propios y baselines, mejor BA logística y limitaciones por falsas alertas, identidad provisional y cobertura asumida. La ampliación financiera añade una segunda pregunta empresarial verificable, con identidad operacional, 16 características y separación entre eventos de pago y estados actuales.

La evaluación financiera mantuvo una población principal de 1\,284 filas y una comparación secundaria de 1\,276 por disponibilidad del cuarto lag histórico. Esta distinción evita que una limitación del baseline cambie la evaluación de clasificadores. Se registraron las decisiones antes de abrir la reserva y se conservaron pesos, sesgos, umbral y preprocesador. La revisión de convergencia y PCA aportó interpretación, no entrenamiento o ajuste adicional del resultado final.

La reutilización segura corresponde a la matemática y a las interfaces verificadas, no a copiar etiquetas, presupuestos, pesos o supuestos entre procesos. Los experimentos muestran distintos compromisos de sensibilidad y especificidad. Sin costos y contexto operacional, no se justifica elegir universalmente una neurona ni convertir sus predicciones en acciones automáticas.

\textbf{La ampliación a los demás Data Marts sigue siendo una obligación académica pendiente.} Pedidos conserva estudios por línea/documento con bloqueo temporal y un objetivo de retraso descartado; Compras conserva estudios de viabilidad sin dataset ni clasificadores aprobados bajo su contrato actual. Se describen como investigaciones con limitaciones, nunca como modelos implementados. Inventario/Movimientos continúa pendiente, así como una revisión final de procesos bloqueados si se obtiene evidencia nueva. Estas actividades no se ejecutaron para ampliar el informe.

La entrega documenta dos aplicaciones y no cierra el proyecto completo. Continúan pendientes la certificación operacional de conservación/cierres, la validación de costos FP/FN, información empresarial adicional, evidencia sobre cuentas nuevas y la identificación inequívoca del ciclo de entrega. Se requiere una evaluación independiente para cualquier uso real; las reservas ya observadas no pueden presentarse de nuevo como test nunca visto.
